\begin{exposition}
The difference quotient definition identifies the derivative as a limiting
slope. Zorich's approximation viewpoint suggests an equivalent reading: near
$a$, a differentiable function is a constant term plus a linear main term plus
an error smaller than the increment. In symbols,
\[
  f(a+h)=f(a)+Lh+r(h),
  \qquad r(h)=o(h)\quad (h\to 0).
\]
The derivative is the coefficient $L$ of the best first-order approximation.
This puts differentiation in the same error-propagation lineage as numerical
approximation and sequence-limit algebra.
\end{exposition}

\begin{definitionbox}{Definition (First-Order Approximation at a Point)}
\begin{definition}[First-Order Approximation at a Point]
\label{def:first-order-approximation-at-a-point}
Let $f:A\to\mathbb{R}$, let $a\in A$, and suppose $a+h\in A$ for all
sufficiently small $h$. A function of the form
\[
  h\longmapsto f(a)+Lh
\]
is called a \textbf{first-order approximation} to $f$ at $a$ with coefficient
$L$.
\end{definition}
\end{definitionbox}

\begin{remark*}[Standard quantified statement]
\[
\operatorname{HasLinearApproximation}(f,a,L)
\quad\text{with approximating map}\quad
h\mapsto f(a)+Lh.
\]
\end{remark*}

\begin{remark*}[Interpretation]
A first-order approximation keeps the value $f(a)$ and adds a linear response
$Lh$ to a small increment $h$. The coefficient $L$ is the candidate derivative
read as the best linear part of the local change in $f$.
\end{remark*}

\begin{dependencies}
\begin{itemize}
  \item \hyperref[def:function]{Function}
\end{itemize}
\end{dependencies}

\begin{definitionbox}{Definition (Differentiability by Linear Approximation)}
\begin{definition}[Differentiability by Linear Approximation]
\label{def:differentiability-linear-approximation}
Let $f:A\to\mathbb{R}$ and let $a\in A$ be an interior point of $A$. We say
that $f$ is \textbf{differentiable at $a$ by linear approximation} with
coefficient $L$ if there is a remainder function $r$ such that
\[
  f(a+h)=f(a)+Lh+r(h)
\]
for all sufficiently small $h$ with $a+h\in A$, and
\[
  r(h)=o(h)\qquad (h\to 0).
\]
\end{definition}
\end{definitionbox}

\begin{remark*}[Standard quantified statement]
\[
\operatorname{IsDifferentiable}(f,a)
\Longleftrightarrow
\exists L\in\mathbb{R}\;\exists r\;
\bigl(f(a+h)=f(a)+Lh+r(h)\ \text{and}\ r(h)=o(h)\bigr).
\]
\end{remark*}

\begin{remark*}[Predicate reading]
The predicate $\operatorname{IsDifferentiable}(f,a)$ can be read as the
existence of a linear coefficient $L$ and a remainder $r$ whose size is
negligible compared with the increment $h$.
\end{remark*}

\begin{remark*}[Interpretation]
The condition $r(h)=o(h)$ is the quantified assertion
\[
  \forall\varepsilon>0\;\exists\delta>0\;\forall h\,
  \bigl(0<|h|<\delta
  \Rightarrow |r(h)|\leq \varepsilon |h|\bigr).
\]
Thus the error is not merely small; it is small relative to the input
increment.
\end{remark*}

\begin{dependencies}
\begin{itemize}
  \item \hyperref[def:first-order-approximation-at-a-point]{First-Order Approximation at a Point}
  \item \hyperref[def:increment-little-o]{Increment Little-o}
\end{itemize}
\end{dependencies}

\begin{theorembox}{Theorem (Derivative and Linear Approximation)}
\begin{theorem}[Derivative and Linear Approximation]
\label{thm:derivative-linear-approximation-equivalence}
Let $f:A\to\mathbb{R}$ and let $a\in A$ be an interior point of $A$. Then
$f$ is differentiable at $a$ with derivative $L$ if and only if
\[
  f(a+h)=f(a)+Lh+o(h)
  \qquad (h\to 0).
\]

\smallskip\noindent
\hyperref[prf:derivative-linear-approximation-equivalence]{\textit{Go to proof.}}
\end{theorem}
\end{theorembox}

\begin{remark*}[Standard quantified statement]
\[
\operatorname{Derivative}(f,a,L)
\Longleftrightarrow
f(a+h)=f(a)+Lh+o(h)\quad(h\to 0).
\]
\end{remark*}

\begin{remark*}[Interpretation]
Subtract $f(a)+Lh$ and divide by $h$:
\[
  \frac{f(a+h)-f(a)-Lh}{h}
  =
  \frac{f(a+h)-f(a)}{h}-L.
\]
The right side tends to zero exactly when the difference quotient tends to
$L$.
\end{remark*}

\begin{dependencies}
\begin{itemize}
  \item \hyperref[def:derivative-at-a-point]{Derivative at a Point}
  \item \hyperref[prop:derivative-h-form-equivalence]{Derivative: $h$-Form Equivalence}
  \item \hyperref[def:differentiability-linear-approximation]{Differentiability by Linear Approximation}
  \item \hyperref[def:increment-little-o]{Increment Little-o}
\end{itemize}
\end{dependencies}

\begin{propositionbox}{Proposition (Uniqueness of the Linear Coefficient)}
\begin{proposition}[Uniqueness of the Linear Coefficient]
\label{prop:linear-approximation-unique}
Let $f:A\to\mathbb{R}$ and let $a\in A$ be an interior point of $A$. If
\[
  f(a+h)=f(a)+Lh+o(h)
  \qquad\text{and}\qquad
  f(a+h)=f(a)+Mh+o(h),
\]
then $L=M$.

\smallskip\noindent
\hyperref[prf:linear-approximation-unique]{\textit{Go to proof.}}
\end{proposition}
\end{propositionbox}

\begin{remark*}[Standard quantified statement]
\[
\bigl(f(a+h)=f(a)+Lh+o(h)\ \land\
f(a+h)=f(a)+Mh+o(h)\bigr)
\Longrightarrow L=M.
\]
\end{remark*}

\begin{remark*}[Interpretation]
The linear coefficient in a first-order approximation is unique. Two different
linear coefficients would leave a nonzero linear term after subtraction, and a
nonzero linear term cannot be absorbed into an $o(h)$ error.
\end{remark*}

\begin{dependencies}
\begin{itemize}
  \item \hyperref[thm:derivative-linear-approximation-equivalence]{Derivative and Linear Approximation}
  \item \hyperref[thm:uniqueness-of-the-derivative]{Uniqueness of the Derivative}
\end{itemize}
\end{dependencies}

\begin{exposition}
Let $f(x)=x^2$. At $a$,
\[
  f(a+h)
  =
  (a+h)^2
  =
  a^2+2ah+h^2.
\]
Since $h^2=o(h)$, the derivative is the coefficient of $h$:
\[
  f'(a)=2a.
\]
\end{exposition}

\begin{exposition}
The statement
\[
  f(a+h)=f(a)+f'(a)h+o(h)
\]
is the first Taylor formula with Peano remainder. Taylor expansion extends
the same idea to higher powers of $h$: keep more main terms and demand a
smaller remainder.
\end{exposition}
